\sectionlink{はじめに}




自然についての注目すべき事実は、興味深い現象が、エネルギー・長さ・時間の広い範囲のスケールにわたって生じることである。この一世紀だけでも、私たちは研究の対象を、原子の大きさの10億分の1ほどの微視的スケール（$10^4$ GeV）から、観測可能な宇宙の大きさに及ぶ巨視的スケール（$10^{-42}$ GeV）まで広げてきた。この巨大なエネルギー範囲こそが、物理科学の成功の鍵である。もしすべての物理現象が同じスケールで起きていたなら、それらを一度にすべて理解しなければならなかっただろう。物理学の進歩は不可能だったはずである！　実際には、このスケールの分離を利用して、関心のある物理過程の典型的エネルギー$E$よりはるかに大きい（小さい）スケールを、第一近似としてすべて無限大（ゼロ）に置くことができる。この極限では、$E$と大きく異なるスケールの物理は無関係となり、無視できる。必要なら、その影響を摂動論で再び取り入れればよい。この基本戦略は、今日に至るまで、あらゆる分野の物理学者によって、多かれ少なかれ暗黙のうちに採用されている。

有効場の理論（EFT）は、上で述べた戦略を明示的に実装し、精密で定量的な枠組みにするための模型構築の道具である。素粒子物理学者はエネルギースケールで考えることを好むが、EFTは、長さ、時間、速度、運動量、角運動量など、あらゆる種類の次元をもつ量の間の階層性を利用して構築できる。実際、適切に定式化すれば、次のようにまで主張できる。
\vspace{.3cm}
%
\begin{quotation}
\noindent \hspace{3cm} {\it すべての物理理論は有効理論である。}
\end{quotation}
%
\vspace{.3cm}

以下では、EFTの主要な道具のいくつかを、さまざまな状況でどのように働くかを示しながら紹介する。本講義ノートは、高エネルギー物理学の基礎知識をもつ大学院生や若手の博士研究員を対象としている。実際、私が大学院生だったときに読みたかったと思うノートを書くことを目指した。したがって、内容の大部分はEFTを使い慣れた読者には既知だろうが、そのような読者にも、ところどころに有益な点を見いだしてもらえれば幸いである。


\subsectionlink{EFTの主要な構成要素}


本講義では、Wilsonによる繰り込み群（RG）へのアプローチの議論から始めるような、歴史的な順序には従わない。主要な原理は非常に簡単に述べられるので、ただちに本題に入り、EFTについての現在の理解を手短にまとめる。EFTは有効作用によって定義され、その有効作用は次の三つの要素によって完全に指定される。
%
\begin{enumerate}
\item {\it 自由度。} EFTを構築する際の第一歩は、関心のある物理系を記述するうえで、どの自由度が重要かを見極めることである。これらが有効作用に現れる変数になる。ここでの鍵は「重要な」という言葉である。追加の構造を導入すれば、どんな現象の記述もいくらでも複雑にできるが、その見返りはすぐに小さくなる。逆に、最も簡潔な記述を目指すこともできるが、Albert Einsteinの言葉を借りれば、{\it 「すべては可能な限り単純にすべきだが、それ以上単純にしてはならない」}。Nambu-Goldstoneモードの場合のように、用いるべき自由度が対称性の考察から示唆されることもある。しかし、多くの場合、自由度の選択は独立した入力であり、科学であると同時に技でもある。
\item {\it 対称性。} EFTを構築する第二段階は、有効作用の形、ひいては系の力学を制約する対称性を特定することである。対称性には、大域的なもの、ゲージ化されたもの、偶発的なもの、自発的に破れたもの、量子異常をもつもの、近似的なもの、縮約されたものなど、さまざまな種類がある。これらすべての例を、順を追って見ていく。ここでは、系の対称性と両立する項は、原理的にはすべて有効作用に含めるべきだとだけ述べておく。その結果、有効作用は一般に無限個の項を含む。
\item {\it 展開パラメータ。} 無限個の項をもつ作用を扱う鍵は、すべてのEFTが一つ以上の展開パラメータを備えているという事実にある。これらは、無視することにした物理が、残すことにした自由度に及ぼしうる影響を制御する小さな量である。例えば素粒子物理学では、これらの展開パラメータはしばしばエネルギースケールの比$E / \Lambda$であり、$E$は関心のある過程の特徴的なエネルギースケール、$\Lambda$は無視する紫外（UV）物理の典型的なエネルギースケールである。他の例としては、速度の比（例えば非相対論的極限での$v/c$）、角運動量の比（半古典的展開の場合）、あるいは小さな無次元結合定数が挙げられる。観測量は、これらの小さなパラメータについての級数として摂動論で計算する。この戦略が機能するには、明示的なパワーカウンティングの方式をもつことが不可欠である。すなわち、有効作用の各項に、展開パラメータについての明確な次数を割り当てられなければならない。これによって、摂動論の任意の次数で寄与する項が有限個に限られ、また、求める精度に応じて作用に残す項をあらかじめ決められることが保証される。
\end{enumerate}
%

これら三つの要素を述べることは比較的容易だが、EFTにそれらを正しく実装するには、微妙な注意が必要になることがある。それを実際にどう実現するかは、扱う問題に依存することが多い。この半世紀の間に豊富なノウハウが蓄積されてきた。本講義ノートの目的は、その主要な手法を紹介することである。



\subsectionlink{EFTの主な利点}


より根本的な理論が存在しないため、物理系のEFTによる記述に頼らざるをえない場合がある。また、そのような理論が存在しても、例えば強結合であるために、厳密な予言はおろか近似的な予言を引き出すことすら難しい場合もある。このような状況では、EFTは同じ物理について、より扱いやすい弱結合の記述を与えうる。さらに、必要に迫られてEFTを使う場合でなくても、EFTは便利な枠組みである。EFTのアプローチの主な利点は、次のようにまとめられる。
%
\begin{itemize}
\item EFTは、重要な自由度に焦点を絞り、扱う問題にとって重要でない側面を最初から無視することで、計算を大幅に簡単にする。
\item 重要な自由度を取り出すと、それまでは隠れていた新たな対称性が明らかになることがある。多くの場合、こうした対称性を用いると、ほとんど計算せずに、あるいはまったく計算せずに、一般的な結論を導ける。
\item EFTを用いると、計算は本質的に二つの部分に分離される。一つは残しておく自由度を含む部分（実際のEFTの計算）であり、もう一つは無視する物理に依存する部分（マッチング計算）である。この「モジュール化された」アプローチによって、同じような計算を何度も繰り返すことを避けられる。
\item 異なるスケールをいくつも含む問題では、それらの比の対数に依存する観測量が現れることがある。スケールが大きく離れていると、対数は大きくなり、摂動計算の精度を悪化させうる。EFTは一度に一つのスケールに注目することで、RGによる走りを使って、これらの大きな対数を足し上げられる。
\item EFTには、無視した物理の系統的なパラメータ表示が含まれる。その物理が未知であったり、十分に理解されていなかったりする場合でも、EFTの予言を実験と比較することで、模型に依存しない制約を課せる。
\item EFTの観測量は、一つ以上の展開パラメータに関する級数である。摂動論をどこまで進めるかによって、EFTの結果の精度が決まる。
\item これに関連して、EFTは、その予言をいつ信頼するのをやめるべきかまで教えてくれる。展開パラメータが$\mathcal{O}(1)$になる領域に入ると、摂動論は機能しなくなる。これは、それまで使っていた自由度がもはや重要な自由度ではなくなり、別のEFTを考える必要があることを意味する。
\end{itemize}
%
ここまでの指摘の中には、現時点では少しわかりにくく感じられるものもあるかもしれない。その場合は、このまま読み進め、EFTの例をいくつか見た後で、この小節に戻ってほしい。 \\

\noindent {\it 規約：} 特に断らない限り、文献~\cite{Srednicki:2007ab}の規約に従う。とりわけ、計量には「ほとんどプラス」の符号規約を用い、$\hbar = c = 1$となる単位系で議論する。また、電磁気学にはGauss単位系を用いる。同じ添字の組は、一方が上付き、他方が下付きであれば、常に縮約されているものと理解する。本講義全体を通じて、導入する重要な概念には\high{下線を引く}。文献は、追加の背景知識や詳細を提供するためにのみ引用し、業績の帰属を示す目的では引用しない。そのため、ご自身の貢献が適切に認められていないと感じられる方々には、あらかじめお詫びしておく。場の量子論については既知とし、標準的な内容については文献~\cite{Srednicki:2007ab}を参照してほしい。


\newpage


\sectionlink{有効場の理論のトイモデル} \label{sec: toy model}



有効場の理論に関する基本的な考え方と手法のいくつかを導入するため、まず次のトイモデルを考えよう。
%
\begin{align} \label{full theory}
	S = \int d^4 x \left\{ - \tfrac{1}{2} \d_\mu \phi \, \d^\mu \phi - \tfrac{1}{2} M^2 \phi^2 + \bar \Psi (i \gamma^\mu \d_\mu - m ) \Psi +  i y \phi \bar  \Psi \gamma^5 \Psi - \tfrac{1}{4!} \lambda \phi^4 \right\} .
\end{align}
%
この作用は、$\phi$が擬スカラー、すなわちパリティ変換の下で$\phi \to - \phi$と変換する場であれば、パリティ不変である。$\phi$の3次の自己相互作用は不変ではないので、省いてある。摂動論を使えるように、$y, \lambda \ll 1$であるパラメータ空間の領域を考える。さらに、スカラーがフェルミオンよりはるかに重い、すなわち$M \gg m$であると仮定する。



ここで、二つのフェルミオンの弾性散乱のような単純な過程を記述したいとしよう。摂動論の最低次、すなわちツリーレベルでは、この相互作用は一つのスカラーの交換によって媒介される。対応するFeynman図を図\ref{fig fermion scattering full}に示す。散乱振幅は、標準的なFeynman則（例えば文献~\cite{Srednicki:2007ab}を参照）を用いて次のように求められる。
%
\begin{align}
	i \mathcal{M} = (- y)^2 \left\{ \bar u_{\lambda_3, p_3} \gamma^5 u_{\lambda_1, p_1} \frac{-i}{(p_3-p_1)^2 +M^2} \,  \bar u_{\lambda_4, p_4} \gamma^5 u_{\lambda_2, p_2} \, - \, (3 \leftrightarrow 4) \right\} . 
\end{align}
%
この過程に関与するフェルミオンの典型的なエネルギー$E$が
%
\begin{align} \label{energy range}
	m \ll E \ll M , 
\end{align}
%
を満たすなら、スカラーの伝播関数を運動量のべきで展開でき、最低次の振幅は
%
\begin{align} \label{4pt amplitude EFT}
	i \mathcal{M} \simeq - i \, \frac{y^2}{M^2} \left[ \bar u_{\lambda_3, p_3} \gamma^5 u_{\lambda_1, p_1}  \bar u_{\lambda_4, p_4} \gamma^5 u_{\lambda_2, p_2} \, - \, (3 \leftrightarrow 4) \right] ,
\end{align}
%
となる。ただし、$\mathcal{O}(E^2 /M^2)$の補正を除いている。興味深いことに、この結果は、次の作用で記述される\high{有効場の理論}（EFT）から出発しても得られる。
%
\begin{align} \label{4pt tree effective theory}
	S_{\rm eff} = \int d^4 x \left\{ \bar \Psi (i \gamma^\mu \d_\mu - m ) \Psi - \frac{y^2}{2 M^2} (\bar  \Psi \gamma^5 \Psi)^2 \right\} 
\end{align}
%
この作用には重いスカラー場がまったく現れない。これは直観にもかなっている。式\eqref{energy range}のエネルギー範囲では、エネルギー保存則によって重いスカラーの生成が禁止されるため、初期状態にスカラーがなければ、起きていることをフェルミオンだけで記述できるはずである。この直観は、まだ発見していない自然の側面を気にかけずに、あらゆる種類の物理現象を模型化してきた成功の歴史に基づいている。このトイモデルの文脈では、$M$を超えるエネルギーに到達できない物理学者でも、実験結果を記述する理論を構築できるはずだ、ということである。

%
\begin{figure}
	\begin{center}
	\begin{fmfgraph*}(90,50) 
	\fmfleft{i1,i2} \fmfright{o1,o2} 
	\fmflabel{$p_1$}{i2}\fmflabel{$p_3$}{o2}
	\fmflabel{$p_2$}{i1}\fmflabel{$p_4$}{o1}
	\fmf{fermion}{i1,v1,o1} \fmf{fermion}{i2,v2,o2} 
	\fmffreeze
	\fmf{dashes}{v1,v2}
	\end{fmfgraph*}
	%
	\hspace{2.5cm}
	%
	\begin{fmfgraph*}(90,50) 
	\fmfleft{i1,i2} \fmfright{o1,o2} 
	\fmflabel{$p_1$}{i2}\fmflabel{$p_4$}{o2}
	\fmflabel{$p_2$}{i1}\fmflabel{$p_3$}{o1}
	\fmf{fermion}{i1,v1,o1} \fmf{fermion}{i2,v2,o2} 
	\fmffreeze
	\fmf{dashes}{v1,v2}
	\end{fmfgraph*}
	%
	\vspace{.3cm}
\caption{\small ツリーレベルにおける$2 \to 2$フェルミオン散乱。}
		\label{fig fermion scattering full}
		\end{center}
\end{figure}
%


重要なのは、$E^2/M^2$が散乱振幅$\mathcal M$の測定精度より小さい限り、理論\eqref{full theory}と\eqref{4pt tree effective theory}は、実質的に区別できない予言を与えることである。したがって、Occamの剃刀に従えば、考えている散乱過程について、同じ精度でより簡単な説明を与えるEFT\eqref{4pt tree effective theory}を使えばよい。比$E/M$は、EFTが完全な理論をどれほどよく近似するかを制御する。これが展開パラメータである。エネルギーが$E > M$になると、EFTは有用ではなくなり、完全な理論に戻る必要がある。

有効作用\eqref{4pt tree effective theory}の注目すべき特徴は、その\high{局所性}である。$S_{\rm eff}$は、同じ点で評価した$\Psi$とその微分の積だけを含む。局所性は、$\phi$がフェルミオン間に媒介する相互作用が、摂動論の最低次では次のYukawaポテンシャルで記述されることから従う。
%
\begin{align}
	V(r) = - \frac{y^2}{4 \pi} \frac{e^{-M r}}{r} \ .
\end{align}
%
考えている過程では、フェルミオンの間で交換される運動量は$q \sim E \ll M$である。したがって、Yukawaポテンシャルを探る分解能は$ \Delta r \sim 1/q \gg 1/M$となり、ポテンシャルの到達距離よりはるかに大きい。この極限では、$V(r)$は本質的に接触相互作用に帰着する。フェルミオンは、不確定性原理によって課される限界の範囲内で、同じ点にいる場合にのみ互いに相互作用できる。




\subsectionlink{ツリーレベル振幅のマッチング}



有効作用\eqref{4pt tree effective theory}は、理論\eqref{full theory}から得られるフェルミオンのツリーレベル4点振幅を、$E/M$の最低次で正しく再現する。そこで、この有効理論のアプローチをどこまで進められるかを問うのは自然である。4点関数について$E/M$の高次補正も再現できるだろうか。より多点の関数はどうか。ループ図はどうか。本小節では最初の二つの問いを扱い、三つ目は\ref{sec: matching at loop level}小節まで後回しにする。

前節の議論によれば、式\eqref{4pt tree effective theory}のEFTには、少数の項だけでなく、すべての対称性を保ちながら$\Psi$とその微分から構成できるあらゆる項を含めるべきである。各項は、先験的には任意の結合定数、すなわち\high{Wilson係数}を掛けた形で有効Lagrangianに入る。これらの係数が無限個あることは、十分に低いエネルギーで単一のDirac場$\Psi$を含み、同じ対称性をもつ理論が、原理的には無限に存在することを反映している。関心のある特定の理論、ここではトイモデル\eqref{full theory}の低エネルギー物理を再現するために、Wilson係数がとるべき正確な値を求める手続きを\high{マッチング}という。実際のマッチング計算は、通常、個々の振幅のレベルで、ループ展開の必要な次数まで行う。これは、前の小節でツリーレベルの4点関数に対して行ったことと同様である。ただし、この手続きが生成汎関数のレベルで直接どのように理解できるかを議論するのも有益である。

散乱振幅は、LSZ簡約公式~\cite{Srednicki:2007ab}を使って、場の時間順序相関関数から取り出せる。一方、任意の時間順序相関関数は、生成汎関数
%
\begin{align} \label{toy model generating functional}
	Z [ J, \bar \eta, \eta] = \int D \phi D \bar \Psi D \Psi \, e^{i S [ \phi, \Psi] + i \int (J \phi + \bar \eta \Psi + \bar \Psi \eta)} ,
\end{align}
%
を外部電流$J, \bar \eta$と$\eta$で適切な回数だけ微分し、その後で電流をゼロに置くことで得られる。フェルミオンが関与する散乱過程にだけ関心があるなら、必要なのは$\Psi$と$\bar \Psi$の相関関数だけである。したがって、式\eqref{toy model generating functional}では最初から$J=0$と置いてよい。その場合、経路積分は次のように分解できる。
%
\begin{align}
	Z [ J=0, \bar \eta, \eta] &= \int  D \bar \Psi D \Psi \, e^{i /\hbar\int (\bar \eta \Psi + \bar \Psi \eta)} \int D \phi \, e^{i S [ \phi, \Psi]/\hbar} \nonumber \\ 
	&\equiv \int  D \bar \Psi D \Psi \, e^{i /\hbar\int (\bar \eta \Psi + \bar \Psi \eta)} e^{i S_{\rm eff} [\Psi]/\hbar} , \label{path integral effective action}
\end{align}
%
ここでは、一時的に$\hbar$の因子をすべて復活させた。この観点から見ると、$\Psi$の有効作用はスカラー場を\high{積分消去}することで得られる。$\phi$に関する積分を厳密に行うと、外線がフェルミオンであるFeynman図に含まれるスカラー内線の寄与が取り込まれる。その結果、有効作用は$S_{\rm eff} = S_{\rm eff}^{(0)} + \hbar S_{\rm eff}^{(1)} + \mathcal{O} (\hbar^2)$という形で$\hbar$のべきに展開でき、$\hbar$に比例する部分はスカラーのループに由来する。この展開の最低次の項$S_{\rm eff}^{(0)}$は、完全な理論\eqref{full theory}におけるフェルミオンのすべてのツリーレベル振幅を再現するEFTを記述する。次の次数の項$S_{\rm eff}^{(1)}$はすべての1ループ補正を再現し、以下も同様である。

ツリーレベルの有効作用$S_{\rm eff}^{(0)}$は、$\phi$についての汎関数積分に鞍点近似を適用して得られるため、とりわけ簡単に計算できる。すなわち、
%
\begin{align} \label{saddle point tree level effective action}
e^{i S_{\rm eff}^{(0)}[\Psi]/\hbar}  = e^{i S [ \hat \phi, \Psi]/\hbar}, \qquad \quad \text{ここで$\hat \phi[\Psi]$は次式の解：} \qquad \quad \frac{\delta S[ \hat \phi, \Psi]}{\delta \hat \phi} = 0 \ .
\end{align}
% 
解$\hat \phi$は、EFTの展開パラメータ$E/M$について摂動的に計算できる。この計算では、場$\Psi$は外場の役割を果たす。作用\eqref{full theory}を$\phi$について変分すると、運動方程式
%
\begin{align}
	(\square -M^2)\phi + y \mathcal{J} - \frac{\lambda}{3!} \phi^3 = 0 ,
\end{align}
%
を得る。ここでは$\mathcal{J} \equiv \bar i \Psi \gamma^5 \Psi$を定義して記法を簡略化した。$E/M$展開の系統的な扱いについての詳しい議論は後回しにし、解を$M$の逆べきで展開することで、この方程式を摂動的に解く。
%
\begin{align}
	\hat \phi = \frac{y}{M^2} \left(1 + \frac{\square}{M^2} + \frac{\square^2}{M^4} + \frac{\square^3}{M^6} \right) \mathcal{J} -\frac{\lambda y^3}{6 M^8} \mathcal{J}^3 +  \cdots \ .
\end{align}
%
この解をトイモデルの作用に代入し、式\eqref{saddle point tree level effective action}を用いると、
%
\begin{align} \label{toy tree level effective action}
	S_{\rm eff}^{(0)} = \int d^4 x \left\{ \bar \Psi (i \gamma^\mu \d_\mu - m ) \Psi + \frac{y^2}{2 M^2} \mathcal{J} \left(1 + \frac{\square}{M^2} + \frac{\square^2}{M^4} + \frac{\square^3}{M^6} \right) \mathcal{J} - \frac{\lambda y^4}{4!M^8}\mathcal{J}^4 + \cdots \right\}.
\end{align}
%
を得る。これは有効作用\eqref{4pt tree effective theory}を改良したもので、4点振幅に限らず、すべてのツリーレベル振幅を取り込んでいる。この作用が無限個の項、実際には対称性と両立するすべての項を含むことは明らかである。各項の前に付くWilson係数は、完全な理論のパラメータ$M, y, \lambda$の非常に特定された関数である。この特定の結合定数の選び方によって、EFTのツリーレベル散乱振幅は、完全な理論を使って得られるものと厳密に同じになる。こうして、EFTをトイモデルへツリーレベルで正しくマッチングできた。

マッチングの手続きを明示的に実行できたのは、（1）$E >M$のエネルギーにおける完全な理論を知っていたことと、（2）その理論が弱結合だったことによる。とりわけ弱結合の仮定のおかげで、重要な自由度をLagrangianから直接読み取り、低エネルギーでの展開パラメータが$E /M$であることを把握できた。しかし、上の二条件のいずれかが満たされないことも多い。その場合、低エネルギーで重要な対称性と自由度を見極めるには、実験、卓越した理論的直観、あるいはその両方に頼る必要がある。

QCDのように、完全な理論は既知で強結合だが、低エネルギー自由度は弱結合である場合には、マッチングを数値的に実行できる。一方、標準模型の場合のように完全な理論が未知でも、任意のWilson係数をもつ有効作用を書くことはできる。これらの係数を実験的に測定すること、あるいはその値に上限を与えるだけでも、背後にある完全な理論について重要な手がかりが得られる。このように、EFTは、関心のある物理を記述し、無視する物理を系統的にパラメータ化するための有用な道具であり続ける。

最後に一点述べておく。マッチングの最終的な目的は、例えばオンシェルの散乱振幅のような物理的観測量について、EFTと完全な理論が厳密に同じ予言をするようにすることである。それでも、生成汎関数を用いたアプローチが示すように、Wilson係数を読み取るために、必ずしも観測量を比較する必要はない。より計算しやすい非物理的な量（例えば、外部運動量をゼロにしたオフシェル振幅）を比較してマッチングを行い、そのようにマッチングしたEFTを使って物理的な予言を導いてもよい。
 


\subsectionlink{パワーカウンティングと強結合スケール}



ここまでで、無限個の項を含む作用に少し不安を覚えた読者もいるかもしれない。Wilson係数を任意のパラメータとして扱う場合は、なおさらである。物理を学ぶ過程で出会う数多くの模型は、どれも少数のパラメータを含む有限のLagrangianで記述される。無限のLagrangianを、いったいどのように扱えばよいのだろうか。すべてのWilson係数を決めるには無限回の測定が必要に思えるが、そのような理論がどうして有用であり、まして予言能力をもちうるのだろうか。これから議論するように、答えはEFTを定義する第三の性質、すなわち展開パラメータにある。

式\eqref{4pt amplitude EFT}の4点振幅は、$E/M$のべき展開の最低次の項だった。実際、EFTのすべての物理的予言は、$E/M$の級数として表されるべきである。この展開をどこまで進めるべきかは、測定可能な精度に依存する。$(E/M)^n$の次数の相対補正が測定できないなら、それを含めても意味はない。EFTは非常に実用主義的な営みである。

少し考えれば、ツリーレベルの有効Lagrangian\eqref{toy tree level effective action}で、$M$の2乗より高いべきによって抑制されるすべての項は、より多点の振幅を生むか、$E/M$の高次で4点関数に寄与することがわかる。式\eqref{4pt amplitude EFT}に示した相互作用だけが、4点関数に$1/M^2$の寄与を与えられる。したがって、ある実験が$E/M$の最低次の4点関数にしか感度をもたないなら、測定すべきパラメータは一つだけである。

より一般に、有限の精度をもつ任意の低エネルギー実験の結果が有限個のパラメータで記述される限り、EFTは有用な枠組みである。その場合、有限回の測定だけでこれらのパラメータの値を決めることができ、実用上必要な範囲でEFTを指定できる。EFTは予言能力をもつ枠組みでもある。例えば、あるエネルギーでいくつかの散乱振幅を測定してパラメータを固定すれば、他の観測量についての予言（例えば、同じ振幅の異なるエネルギーでの値の予言）に使えるからである。

実際には、どんな実験についても、有効Lagrangianのどの項を残し、どの項を無視してよいかを、最初の段階で判断できなければならない。各項が観測量に$E/M$の何次で寄与するかをあらかじめ判定できれば、それが可能になる。有効作用の各項が展開パラメータに対してどうスケールするかを求める手続きを、\high{パワーカウンティング}という。このトイEFTで実際にどう働くかを見てみよう。

式\eqref{energy range}のエネルギー範囲では、散乱過程に関与する各フェルミオンの4元運動量は、エネルギーに対して$p^\mu \sim E$とスケールする。この領域では、$\Psi$に作用する各微分も同じようにスケールし、$\d_\mu \sim E$となる。不確定性原理を用いると、$x_\mu \sim 1/E$および$d^4 x \sim 1/E^4$も導ける。最後に、弱結合の有効理論では、運動項が$E/M$展開の最低次の項であると仮定するのが妥当だろう。この要請によって、$\Psi$のスケーリングが決まる。
%
\begin{align} \label	{scaling Psi toy model}
\int d^4 x \, \bar \Psi i \gamma^\mu \d_\mu \Psi \sim 1 \qquad\quad  \Longrightarrow \qquad\quad  \Psi \sim E^{3/2} .
\end{align}
%

これらのスケーリング則を使うと、Lagrangianの各演算子に$E/M$の明確なべきを対応させられる。例えば、式\eqref{toy tree level effective action}の最後の項は、次のようにスケールする。
%
\begin{align}
	\int d^4 x \left\{ - \frac{\lambda y^4}{4!M^8}\mathcal{J}^4 \right\} \sim \left(\frac{E}{M} \right)^8.
\end{align}
%
この単純な例では、マッチングによって得られたWilson係数が$1/M^8$に比例するという事実から、この項が展開パラメータに対してどうスケールするかが、すでに示唆されていた。後で見るように、Lagrangianの演算子のスケーリングを求めることは、しばしばこれほど容易ではない。それでも重要なのは、EFTが明確に定義されるためには、各演算子が展開パラメータの確定したべきとしてスケールしなければならない、という点である。

ここまで導いたスケーリング則によって、Lagrangianの異なる項の相対的な重要性を判断できる。これは、これまで出会ってきたLagrangianが少数の項しか含んでいなかった理由をある程度明らかにする。多くの応用では、展開パラメータが十分に小さいため、最も重要な項だけを考えれば十分なのである。

一方、$E \sim M$になると、すべての相互作用が運動項と同じくらい重要になり、摂動論は機能しなくなる。このため、スケール$M$は通常、EFTの\high{強結合スケール}と呼ばれる。よりくだけた言い方では、$M$をEFTの\high{カットオフ}と呼ぶ人も多い。これは、Wilsonの繰り込み群のアプローチでは、このスケールで遅いモードのループが打ち切られることを示唆している。以下ではこの用語を使わないことにする。というのも、次の小節で議論するように、ループレベルで一貫したパワーカウンティングを行うには、次元正則化のような質量に依存しない正則化法が必要だからである。

次に、スケーリング則を使って、Lagrangianの各項が観測量に与える最低次の寄与をどのように求めるかを議論しよう。散乱振幅については、$n$個のフェルミオンが関与する散乱過程$i \to f$の確率振幅を$n$点関数に結び付けるLSZ簡約公式~\cite{Srednicki:2007ab}を使える。ほぼ質量のないフェルミオンについて、この公式は非常に模式的に次のように書ける。
%
\begin{align} \label{LSZ formula}
	\langle f | i \rangle \equiv (2 \pi)^4 \delta^4 (p_f - p_i) \, i \mathcal{M}_n \sim \left\{ \prod_{j=1}^n \int d^4 x_j \, e^{\pm i p_j \cdot x_j} \slashed{\partial} \, u_j \right\} \langle T \Psi (x_1) \dots \Psi (x_n) \rangle .
\end{align}
% 
この式は、かなり大まかに書いてある。$u_j$の一部は$\bar u_j$、$v_j$、$\bar v_j$かもしれず、$\Psi$の一部は$\bar \Psi$かもしれないからである。さらに、偏極の添字をすべて省略し、各因子の順序にもあまり注意を払っていない。振幅$\mathcal M_n$が$E/M$に対してどうスケールするかを求めるという目的には、こうした細部は重要ではない。

ただし、モード関数$u_j$がエネルギーに対してどうスケールするかを知る必要がある。そのために、自由場の演算子$\Psi$を、生成・消滅演算子で次のように展開できることを思い出そう。
%
\begin{align}
	\Psi (x) = \sum_\lambda \int \frac{ d^3 p}{(2 \pi)^3 2 E_{\vec p}} \left[ b_{\lambda, \vec p} \, u_\lambda (\vec p) e^{i \vec p \cdot x} + d^\dag_{\lambda, \vec p} v_\lambda (\vec p) e^{-i \vec p \cdot x} \right] ,
\end{align}
%
ここで$\lambda$は偏極の添字であり、$[ b_{\lambda, \vec p}, b^\dag_{\lambda', \vec p'} ] = [ d_{\lambda, \vec p}, d^\dag_{\lambda', \vec p'} ] = 2 E_{\vec p} \, (2\pi)^3 \delta (\vec p - \vec p') \delta_{\lambda \lambda'}.$である。これらの交換関係によって生成・消滅演算子のスケーリングが決まり、式\eqref{scaling Psi toy model}を使うと、$u_j \sim E^{1/2}$と結論できる。\footnote{これは、$u_\lambda (\vec p) \bar u_\lambda (\vec p) \sim \slashed{p}$という事実とも整合する。~\cite{Srednicki:2007ab}}

最後に、式\eqref{LSZ formula}の右辺にある$n$点関数は、次式で与えられる。
%
\begin{align} \label{expectation values}
	\langle T \Psi (x_1) \dots \Psi (x_n) \rangle = \frac{\int  D \bar \Psi D \Psi \, \Psi (x_1) \dots \Psi (x_n) e^{i S_{\rm eff} [\Psi]}}{\int  D \bar \Psi D \Psi \, e^{i S_{\rm eff} [\Psi]}} \ .
\end{align}
%
この結果をLSZ公式および本小節で導いたすべてのスケーリング則と組み合わせると、Lagrangianの各項が摂動論で振幅$\mathcal M_n$にどのように寄与するかを見積もれる。例えば、有効作用の中の$n$個の場を含む項$\Delta S$が$n$点関数に与える摂動論の最低次の寄与は、$\Psi (x_1) \dots \Psi (x_n) \Delta S$と同じようにスケールする。

\begin{eBox}
%
{\bf 演習2.1：} 4次の相互作用$\Delta S = \int d^4 x \, \mathcal{O} (x)$がフェルミオンの4点振幅$\mathcal M_4$に与える寄与は、$\mathcal{O} (x)$と同じようにスケールすることを示せ。
%
\end{eBox}
%
\begin{eBox}
%
{\bf 演習2.2：} 式\eqref{toy tree level effective action}の有効理論において、最低次の6点振幅$\mathcal M_6$がエネルギーに対してどうスケールするかを求めよ。
%
\end{eBox}



\subsectionlink{ループレベルでのマッチング} \label{sec: matching at loop level}



ここまで、作用\eqref{toy tree level effective action}で記述されるEFTが、トイモデル\eqref{full theory}を用いて低エネルギーで計算されるすべてのツリーレベル散乱振幅を再現できると論じてきた。次に、マッチングの手続きをループレベルでも正しく実行できることを示す。ツリーレベルの有効作用$S^{(0)}_{\rm eff}$から構成したループ図だけでは、完全な理論のすべての1ループ補正を再現できないことがわかるだろう。もちろん、これは問題ではない。$S^{(0)}_{\rm eff}$は有効作用のツリーレベルの部分にすぎないからである。EFTと完全な理論のループ補正の間に見えるこの不一致は、有効作用の1ループ部分$S^{(1)}_{\rm eff}$によって解消される。実際、$S^{(1)}_{\rm eff}$を計算するには、式\eqref{path integral effective action}の経路積分を計算するよりも、ループ図を比較する方が効率的なことが多い。この方法の利点は、高次のループになるほどいっそう顕著になる。この手続きを説明するため、簡単のためにフェルミオンの2点関数の1ループ補正に注目し、$E/M$展開の最低次で計算する。これで、ループの寄与をマッチングする際に現れる主な注意点を紹介するには十分である。

まず、EFTにおける1ループ補正を考えよう。ここでの目的には、完全な伝播関数ではなく、1粒子既約（1PI）図を計算すれば十分である。$E/M$の最低次では、フェルミオンの2点関数に寄与するFeynman図は一つだけである。
%
\vspace{.9cm}
\begin{align} \label{one loop toy EFT 1}
	%
	\parbox[t][9.5mm][b]{20mm}{\begin{fmfgraph*}(50,60) 
	\fmfleft{i} \fmfright{o} 
	\fmf{fermion}{i,v,o} 
	\fmffreeze
	\fmf{fermion,tension=.5}{v,v}
	\end{fmfgraph*}} = - i \, \frac{y^2}{M^2} \int \frac{d^4 k}{(2 \pi)^4} \, \gamma^5 \, \frac{-i  (- \slashed{k} + m)}{k^2 + m^2 - i \epsilon} \, \gamma^5 = - \, \frac{y^2 m}{M^2} \int \frac{d^4 k}{(2 \pi)^4} \, \frac{1}{k^2 + m^2 - i \epsilon} ,
	%
\end{align}
%
最後の段階では、$(\gamma^5) =1$と、$k$に線形な寄与が対称性によって消えることを使った。この図がUVで2次発散することは明らかである。物理的な考察から、$k$積分をスケール$M$で打ち切って、この発散を正則化したくなるかもしれない。結局のところ、EFTが有効であると想定しているのは、展開パラメータが1より小さいエネルギー$E < M$だけである。さらに、このアプローチは、「速い」Fourierモードを積分消去して、粗視化された観測量の低エネルギー理論を得るという、Wilsonの繰り込み群の見方~\cite{Wilson:1993dy}にも通じる。しかし、これはよい考えではない。そうすると、積分から生じる$M^2$の因子がWilson係数の$1/M^2$と打ち消し合い、UVスケール$M$に依存しない結果になる。これは、4次の自己相互作用は常に$E^2 / M^2$の因子で抑制されるべきだというパワーカウンティングの規則に、直接反する。

この問題に対する最も一般的な解決法は、代わりに次元正則化を用いることである。この方法には複数の利点がある。ゲージ理論ではゲージ不変性を保ち、修正最小減算法（略して$\overline{MS}$法）と組み合わせると繰り込みの手続きを大幅に簡単にし、さらに、これから示すようにパワーカウンティングを損なわない。$d = 4 -\varepsilon$次元に移ると、\footnote{$d=4$からのずれを表すために、丸みを帯びたイプシロン（$\varepsilon$）を使っている。これは、例えば式\eqref{one loop toy EFT 1}の分母に現れる、別の書体のイプシロン（$\epsilon$）と混同してはならない。}式\eqref{one loop toy EFT 1}の右辺は次のようになる。
%
\begin{align} \label{one loop toy EFT 2}
	- \, \frac{y^2 \tilde{\mu}^\varepsilon m}{M^2} \int \frac{d^{4-\varepsilon} k}{(2 \pi)^{4-\varepsilon}} \, \frac{1}{k^2 + m^2 - i \epsilon} =\frac{i y^2}{16 \pi^2} \frac{m^3}{M^2} \left[ \frac{2}{\varepsilon} + 1 + \log\left(\frac{\mu^2}{m^2}\right) + \mathcal{O}(\varepsilon) \right] ,
\end{align}
%
ここでは、$d=4$から離れても$y$が無次元であり続けるように、$y \to y \, \tilde{\mu}^{\varepsilon/2}$と置き換えて、繰り込みスケール$\tilde \mu$を導入した。また、通常どおり、関連するスケール$\mu = \sqrt{4 \pi} e^{-\gamma} \tilde \mu$を導入して記法を簡略化した。ここで$\gamma = 0.55772...$はEuler定数である。ループ積分は今や、被積分関数に現れる唯一のスケール、すなわちフェルミオン質量$m$に比例しなければならない。その結果、1ループ補正は、パワーカウンティングから期待されるとおり、$m^2 /M^2 = m^2 /E^2 \times E^2 / M^2$の因子で抑制される。
%
\begin{eBox}
%
{\bf 演習2.3：} 式\eqref{one loop toy EFT 2}の結果を導け。
%
\end{eBox}
%

次に、完全な理論に注目しよう。EFTでは、フェルミオンの1PI図を計算した。したがって今度は、フェルミオンに関しては1PIだが、重いスカラーに関しては必ずしもそうでない、すべての1ループ図を計算すべきである。このため、1ループで原理的に考慮すべきFeynman図は二つある。一つ目は次の図である。
\vspace{3mm}

%
\begin{align}
	\parbox[t][11.5mm][b]{24mm}{\begin{fmfgraph*}(60,70) 
	\fmfleft{i} 
	\fmfright{o} 
	\fmftop{t}
	\fmf{fermion}{i,v1,o} 
	\fmffreeze
	\fmf{dashes,tension=.7}{v1,v2}
	\fmf{fermion,left,tension=.3}{v2,t,v2} 
	\end{fmfgraph*}} = (-y)^2 \gamma^5 \, \frac{-i}{M^2} \int \frac{d^4 k}{(2 \pi)^4} \,   \mbox{Tr} \left[ \gamma^5 \, \frac{-i  (- \slashed{k} + m)}{k^2 + m^2 - i \epsilon} \right] = 0 .  
\end{align}
%
この図はゼロになる。外部フェルミオン線を取り除いて得られるタッドポール図は、パリティを保つためにゼロでなければならないからである。完全な理論には、もう一つループ補正がある。すなわち、
\vspace{3mm}
%
\begin{align}
	\parbox[t][8mm][b]{30mm}{\begin{fmfgraph*}(80,50) 
	\fmfleft{i} \fmfright{o} 
	\fmf{fermion}{i,v1,v2,o} 
	\fmffreeze
	\fmf{dashes,left}{v1,v2}
	\end{fmfgraph*}} & = (-y)^2 \int \frac{d^4 k}{(2 \pi)^4} \,  \frac{(-i)^2 \gamma^5 (- \slashed{k} - \slashed{p} + m)\gamma^5}{[(k+p)^2 + m^2 - i \epsilon](k^2 + M^2 - i \epsilon)} \label{1 loop full theory 1} \\
	& = - y^2 \int_0^1 dx \int \frac{d^4 q}{(2 \pi)^4} \frac{\slashed{q} + (1-x) \slashed{p} +m}{(q^2 + D)^2} \ ,
\end{align}
%
最後の段階では、分母をまとめるためにFeynmanパラメータ$x$を導入し、$D \equiv x(1-x) p^2 + x m^2 + (1-x)M^2$と定義した。EFTの結果と比較できるように、このループ積分も次元正則化を用いて正則化する。その結果、次式を得る。
%
\begin{align}
& - y^2 \tilde \mu^\varepsilon \int_0^1 dx \int \frac{d^{4-\varepsilon} q}{(2 \pi)^{4-\varepsilon}} \frac{\slashed{q} + (1-x) \slashed{p} +m}{(q^2 + D)^2} = \nonumber \\
& \qquad \qquad \qquad = \frac{i y^2}{16 \pi^2}\left\{ - \frac{1}{\varepsilon} (\slashed{p} + 2 m) + \int_0^1 dx \left[ (1-x) \slashed{p} + m \right] \log (D/\mu^2) + \mathcal{O}(\varepsilon)\right\} \ . \label{1 loop full theory 2}
\end{align}
%
EFTでの計算は$\mathcal{O}(E^2/M^2)$までしか有効でなかったため、比較する前に、式\eqref{1 loop full theory 2}を$m^2/M^2$および$p^2 / M^2$のべきで展開する必要がある。$\mathcal{O} (M^{-2})$までのすべての項を残すと、次の結果が得られる。

%
\begin{align}
	\parbox[t][11.5mm][b]{30mm}{\begin{fmfgraph*}(80,70) 
	\fmfleft{i} \fmfright{o} 
	\fmf{fermion}{i,v1,v2,o} 
	\fmffreeze
	\fmf{dashes,left}{v1,v2}
	\end{fmfgraph*}} & = \frac{i y^2}{16 \pi^2} \bigg\{ - \frac{1}{\varepsilon} (\slashed{p} + 2 m) + \frac{\slashed{p}}{2} \left[ - \frac{1}{2} + \log \left( \frac{M^2}{\mu^2} \right) + \frac{m^2}{M^2} \right] \label{one loop full theory 3} \\
	&  \qquad   + m \left[ -1 + \log \left( \frac{M^2}{\mu^2} \right) + \frac{m^2}{M^2} \log \left( \frac{M^2}{m^2}\right) + \frac{p^2}{M^2}  \right] + \mathcal{O}(\varepsilon, p^3, M^{-4}) \bigg\} . \nonumber 
\end{align}
%

%
\begin{eBox}
%
{\bf 演習2.4：} 式\eqref{one loop full theory 3}の結果を導け。
%
\end{eBox}
%

EFTでの式\eqref{one loop toy EFT 2}と、完全な理論での式\eqref{one loop full theory 3}の両方について、フェルミオンの2点関数の1ループ補正を計算したので、いくつかの結論を導ける。
%
\begin{enumerate}
\item 式\eqref{one loop toy EFT 1}のEFTのループ積分は、式\eqref{1 loop full theory 1}の完全な理論のループ積分において、被積分関数を$1/M$のべきで展開すれば得られる。図のレベルでは、これは、完全な理論の図でスカラーの伝播関数（破線）を一点に縮めると、EFTのFeynman図が得られることに対応する。
\item この関係があるにもかかわらず、式\eqref{one loop toy EFT 2}と\eqref{one loop full theory 3}の最終結果は異なる。これは、ループ積分と級数展開が交換可能な操作ではないことを意味する。特に、$M =\infty$のまわりで非解析的な項が、EFTの結果にはすべて欠けていることに注意しよう。$M$の逆べき展開そのものが解析性の仮定に依存しているので、これはもっともである。
\item $1 /\varepsilon$に比例するUV発散部分も、EFTと完全な理論では異なる。このことは、両理論の質量パラメータとフェルミオン場が異なる異常次元をもつことを意味し、したがって両者を混同してはならない。
\item これに対して、$m=0$のまわりで非解析的な項は同一である。これは常に成り立つ必要があり、実際には整合性の検査に利用できる。次の小節で議論するように、この相殺は、EFTが赤外（IR）とUVのスケールを分離し、大きな対数を制御するうえで不可欠である。
\item 最後に、EFTのループ補正は、完全な理論とは異なり、小さな比$m^2 / M^2$によって自動的に抑制される。例えば式\eqref{one loop toy EFT 2}を参照されたい。このことによって、$E/M$の任意の次数で計算すべきループ図は有限個で済む。
\end{enumerate}
%

以下では、EFTと完全な理論の両方で、$\overline{MS}$法を用いて繰り込みを行う。実際の計算では、これは最終結果から$1/\varepsilon$の因子をすべて除くことに等しい。こうすると、両理論の繰り込まれた1ループ図の不一致は、EFTに次のツリーレベル図を追加することで解消できる。
%
\begin{align}
	& \parbox[t][8mm][b]{30mm}{\begin{fmfgraph*}(78,50) 
	\fmfleft{i} \fmfright{o} 
	\fmf{fermion}{i,v,o} 
	\fmffreeze
	\fmfv{d.sh=circle,d.f=empty,d.si=.2w,l=$\!\!\!\!\! 1$}{v}
\end{fmfgraph*}}  = \left[ \begin{array}{c} \text{1ループ} \\ \text{完全な理論}
\end{array} \right]_{\overline{MS}} - \quad \left[ \begin{array}{c} \text{1ループ} \\ \text{EFT}
	\end{array} \right]_{\overline{MS}} \nonumber \\
	& =  \frac{i y^2}{16 \pi^2} \left\{ - \slashed{p} \left[ \frac{1}{4}-\frac{m^2}{M^2} - \log \left( \frac{M^2}{\mu^2} \right) \right] - m\left(1+\frac{m^2}{M^2} \right) \left[ 1 - \log \left(\frac{M^2}{\mu^2} \right) \right] + \frac{m}{M^2} \, p^2 \right\} . \label{toy matching correction}
\end{align}
%
この図は有効作用の1ループ補正から生じるので、その補正は次の形でなければならない。
%
\begin{align}
	S^{(1)}_{\rm eff} & = \int d^4 x \,  \frac{y^2}{16 \pi^2} \bigg\{ \left[ \frac{1}{4}-\frac{m^2}{M^2} - \log \left( \frac{M^2}{\mu^2} \right) \right] \bar \Psi i \gamma^\mu \partial_\mu \Psi \nonumber  \\
	&   \qquad \qquad \qquad \qquad \qquad - m\left(1+\frac{m^2}{M^2} \right) \left[ 1 - \log \left(\frac{M^2}{\mu^2} \right) \right] \bar \Psi \Psi - \frac{m}{M^2} \bar \Psi \square \Psi + \cdots \bigg\} ,  \label{toy one loop effective action}
\end{align}
% 
ここで省略記号は、$\Psi$について3次以上の項を表す。最初の二項は、ツリーレベルの有効作用\eqref{toy tree level effective action}にある同種の項とまとめられることに注意しよう。一方、最後の項は新しい。これはツリーレベルでは現れなかったが、対称性が許すすべての項が有効作用に含まれるという原則に従って、ループレベルで生成される。

最後に一点述べておく。ここではオフシェル、すなわち任意の$p^2$についてマッチングを行うことを選んだが、オンシェル、すなわち$p^2 = - m^2$を満たす運動量で行ってもよかった。その場合、式\eqref{toy matching correction}の括弧内の最後の項は$-m^3 / M^2$に変わる。したがって、1ループ有効作用\eqref{toy one loop effective action}の最後の項は、$\bar \Psi \Psi$に比例する追加の寄与に置き換わる。この新しい有効作用は、ここで導いた作用と物理的に等価である。それは、ここで得た作用に（摂動的な）場の再定義~\cite{Georgi:1991ch}を施すことで得られる。この操作は$S$行列を変えない~\cite{Weinberg:1995mt}。この場の再定義は、最低次の運動方程式を使って高階微分の項を除くことに等しい。なお、これは\high{一部の演算子が冗長である}ことを示している。すなわち、それらの効果は、他の演算子のWilson係数を変更することで取り込める。オンシェルでマッチングを行うと、冗長な演算子は自動的に除かれる。
 
%
\begin{eBox}
%
{\bf 演習2.5：} オフシェルとオンシェルのマッチングで得られる有効作用を結び付ける、場の再定義を導け。
%
\end{eBox}
%
 
以上では、EFTと完全な理論における2点振幅の1ループ補正を比較することで、$S^{(1)}_{\rm eff}$の2次の部分を取り出した。同様に、両理論のより高次のループ振幅や、より多点の振幅を比較すれば、それぞれ有効作用の高次ループ補正や非線形補正を計算できる。
 


\subsectionlink{大きな対数の足し上げ}




ここまで、EFTは、関心のある観測量にとって重要な自由度に焦点を絞るため、使いやすいと論じてきた。それでも、このトイモデルの具体例では、わずかな項しかもたないLagrangianから、原理的には無限個の項を含む有効Lagrangianへ移ることになった。有限の精度の計算には有限個の項しか寄与しないと論じたものの、なお代償が大きすぎると感じる読者もいるだろう。そうした読者を励ますため、複数のスケールをもつ理論で計算を整理するうえで、EFTが賢明な方法である別の理由を議論しよう。

繰り込みスケール$\mu$の導入は、繰り込みの手続きに必然的に伴うものである。これは先験的には任意のスケールなので、物理的な予言はこのスケールに依存してはならない。それを保証するため、繰り込まれた結合定数は、繰り込み群（RG）方程式に従う形で$\mu$に依存しなければならない。

摂動計算では、最終結果に大きな対数が現れないように$\mu$を選ぶと便利である。例えば、式\eqref{one loop full theory 3}の完全な理論の1ループ補正を考えよう。$\slashed{p}$に比例する部分には、$\frac{y^2}{16 \pi^2} \log (M/\mu)$という形の項が含まれる。同様に、$n$ループ補正には$\left[ \frac{y^2}{16 \pi^2} \log (M/\mu) \right]^n$の次数の項が現れるので、$\log (M /\mu)$が大きいと、摂動論を素朴に適用したときの収束は遅くなる。これは$\mu \sim M$を選べば避けられる。\footnote{そのためには、スケール$M$における結合定数の値を知る必要がある。別のスケールでしか結合定数がわかっていない場合でも、RG方程式を解くことで、$\mu \sim M$での値へ常に外挿できる。}異なるスケールの間に大きな階層がある理論では、対数がUVスケールとIRスケールの比にも依存しうる。例えば式\eqref{one loop full theory 3}の2行目の$\log (M^2 /m^2)$を見てほしい。そのため、大きな対数は避けられないように思えるかもしれない。これから、低エネルギーでEFTを使うことで、この問題を容易に回避できると示す。

前の小節で計算した有効作用をもう一度考えよう。ツリーレベルと1ループの寄与を加えると、2次の部分は次のようになる。
%
\begin{align}
	S^{(0)}_{\rm eff} + S^{(1)}_{\rm eff} & = \int d^4 x \,\left[ \mathcal{Z}(\mu) \bar \Psi i \gamma^\mu \partial_\mu \Psi -  \mathcal{M}(\mu) \bar \Psi \Psi - \mathcal{N}(\mu) \bar \Psi \square \Psi + \cdots \right] ,
\end{align}
%
ここでも、省略記号は作用の非線形な部分を表す。また、次を導入した。
%
\begin{subequations} \label{EFT wilson coefficients}
\begin{align}
	\mathcal{Z}(\mu) &= 1 + \frac{y^2}{16 \pi^2} \left[ \frac{1}{4}-\frac{m^2}{M^2} - \log \left( \frac{M^2}{\mu^2} \right) \right] , \label{Z 1-loop toy} \\
	\mathcal{M}(\mu) &= m \left\{ 1 + \frac{y^2}{16 \pi^2}\left(1+\frac{m^2}{M^2} \right) \left[ 1 - \log \left(\frac{M^2}{\mu^2} \right) \right] \right\} , \label{M 1-loop toy} \\
	\mathcal{N}(\mu) &= \frac{m}{M^2} \ . \label{N 1-loop toy}
\end{align}
\end{subequations}
%
この有効作用は、完全な理論の1ループの2点関数を、$\mathcal{O}(E^2/M^2)$を含む次数まで正確に再現すると論じた。しかし、ここで非常に興味深いことが起きている。

一方で、式\eqref{EFT wilson coefficients}のWilson係数は、$\log (M^2 / \mu^2)$だけに依存する。これは、IRスケール$m$の対数がマッチングで消えていなければ成り立たなかったことに注意しよう。この注目すべき事実は、前の小節でも強調した。したがって、$\mu = M$でマッチングの手続きを行えば、有効作用の大きな対数を完全に避けられる。つまり、EFTのWilson係数の値は、UVスケール$M$で計算すべきなのである。

他方、式\eqref{one loop toy EFT 2}のEFTの1ループ補正は、$\log (m^2 / \mu^2)$だけに依存する。したがって、IRスケール$\mu = m$で1ループ補正を計算すれば、2点関数の大きな対数を避けられる。もちろん、そのためにはIRスケール$m$でのWilson係数の値を知る必要がある。それらは、RG方程式を解き、Wilson係数を、マッチングによって値が得られる$\mu = M$から$\mu = m$まで走らせることで求められる。

要するに、エネルギー$E < M$で完全な理論からEFTへ切り替えることによって、完全な理論に含まれる大きな対数を、本質的に次の二つの部分へ分解している。
%
\begin{align}
	\log \left( \frac{M^2}{m^2} \right) = \log \left( \frac{M^2}{\mu^2} \right) - \log \left( \frac{m^2}{\mu^2} \right) \ . 
\end{align}
%
最初の部分はEFTのWilson係数に吸収され、第二の部分はRG方程式を解くことで「再総和」される。大きな対数を制御するのは、\high{マッチングと走り}の組合せであることを強調しておきたい。完全な理論でRGによる走りを用いるだけでは不十分である。RGが「再総和」できるのは、繰り込みスケール$\mu$に依存する対数だけだからである。




\subsectionlink{繰り込み可能な理論に対する現代的な見方} \label{sec: renormalizable theories}




歴史的に、素粒子物理学のコミュニティでは、繰り込み可能な場の量子論が非常に重視されてきた。実際、式\eqref{full theory}で出発点としたトイモデルは、繰り込み可能であった。この特定の種類の理論には、ループ計算で生じるあらゆるUV発散を、場と有限個の結合定数の再定義に吸収できるという特別な性質がある。当初は、理論が予言能力をもつために、この性質が必要だと考えられていた。繰り込み可能性は、素粒子物理の模型に対する根本的な要請と見なされることが多かったほどである。今日では、EFTは無限個の「繰り込み不可能な」相互作用を含むにもかかわらず、その適用範囲内では同じように予言能力をもつことがわかっている。ループレベルのマッチングについて議論した今なら、繰り込み可能な理論の特別な位置づけを理解できる。

EFTのWilson係数は、UVスケール$M$に二通りの形で依存する。（1）比$m^2/M^2$のべきと、（2）$\log (M^2/\mu^2)$という形の対数である。Wilson係数をスケール$\mu = M$で評価すると、対数はすべて消え、$M$の逆べきが残る。それらは、「繰り込み不可能な」（または\emph{非関連な}）項を抑制し（例えば式\eqref{N 1-loop toy}を参照）、同時に「繰り込み可能な」（または\emph{関連する}）項の係数にも補正を与える（例えば式\eqref{Z 1-loop toy}と\eqref{M 1-loop toy}を参照）。したがって、$m$と$y$を固定して形式的な極限$M \to \infty$をとると、EFTは、繰り込み可能な理論（この場合には自由理論）に帰着する。

より一般に、繰り込み可能な理論は、UVスケールをすべてのIRスケールより無限に大きくし、UV物理が完全に分離する極限で現れる。そのとき、繰り込み可能な理論の裸の結合定数は、$M = \infty$で評価したWilson係数の有限部分集合にすぎない。ループ計算で生じる発散は、今やすべて、裸の結合定数と場の再定義に吸収できる。それ以外の発散は、$1/M \to 0$のべきで抑制されるからである。（ここでは主旨を伝えるために、$M \to \infty$と$\varepsilon \to 0$の極限の順序を、明らかに大まかに扱っている。）



\subsectionlink{自然性} \label{sec: naturalness}



式\eqref{full theory}のトイモデルを、フェルミオンがスカラーよりはるかに重い、すなわち$m \gg M$となる別のパラメータ領域で考えると有益である。この場合も、本節でここまで展開してきた論理に従い、フェルミオンを積分消去することで、エネルギー$E \ll m$で有効なスカラーのEFTを得られる。

ツリーレベルの有効作用$S^{(0)}_{\rm eff}$が、式\eqref{full theory}で$\Psi = 0$とした作用に等しいことは、容易にわかる。1ループ有効作用$S^{(1)}_{\rm eff}$の2次の部分を計算するため、ツリーレベルのEFTと完全な理論から得られる、スカラー伝播関数の1ループ補正を比較する。次元正則化を用いると、EFTでは次の結果が得られる。
\vspace{.7cm} 
%
\begin{align}
%
	\parbox[t][9.5mm][b]{20mm}{\begin{fmfgraph*}(50,60) 
	\fmfleft{i} \fmfright{o} 
	\fmf{dashes}{i,v,o} 
	\fmffreeze
	\fmf{dashes,tension=.5}{v,v}
	\end{fmfgraph*}}	= - \frac{i \lambda \tilde{\mu}^\varepsilon}{2} \int \frac{d^{4-\varepsilon} k}{(2 \pi)^{4-\varepsilon}} \frac{-i}{k^2 + M^2 - i \epsilon} = \frac{i \lambda}{32 \pi^2} M^2 \left[ \frac{2}{\varepsilon} + 1 + \log\left(\frac{\mu^2}{M^2}\right) + \mathcal{O}(\varepsilon) \right] . \label{toy EFT scalar loop diagram}
\end{align}
%

同じ図は完全な理論にも現れるが、それに加えて、ループ内をフェルミオンが流れる第二の図がある。
%
\begin{align}
	\parbox[t][11.5mm][b]{24mm}{\begin{fmfgraph*}(60,70) 
	\fmfleft{i} 
	\fmfright{o} 
	\fmf{dashes}{i,v1} 
	\fmf{dashes}{v2,o} 
	\fmf{fermion,left,tension=.3}{v1,v2,v1} 
	\end{fmfgraph*}} \!\! = \,\, - (- y)^2 \int \frac{d^4 k}{(2 \pi)^4}\text{Tr} \left[ \frac{-i(- \slashed{k} - \slashed{p} + m)}{(k+p)^2 + m^2 - i \epsilon} \gamma^5 \frac{-i(- \slashed{k} +m)}{k^2 + m^2 - i \epsilon}\gamma^5 \right] . 
\end{align}
%
一点注意しておく。行列$\gamma^5$は$d=4$から離れると、適切に定義されない。したがって、$d = 4 - \varepsilon$に移る前に、$(\gamma^5)^2 = 1$および$\gamma^5 \slashed{p} \gamma^5 = - \slashed{p}$を用いてこの行列を除く必要がある。この点に注意して計算し、最終結果を$p^2 / m^2$のべきで展開すると、次式を得る。
%
\begin{align}
	\parbox[t][11.5mm][b]{24mm}{\begin{fmfgraph*}(60,70) 
	\fmfleft{i} 
	\fmfright{o} 
	\fmf{dashes}{i,v1} 
	\fmf{dashes}{v2,o} 
	\fmf{fermion,left,tension=.3}{v1,v2,v1} 
	\end{fmfgraph*}} \!\! = \,\, - \frac{i y^2}{4 \pi^2} \left\{ \frac{1}{\varepsilon} (p^2 + 2 m^2) +\frac{p^2}{2} \log \left( \frac{\mu^2}{m^2} \right) + m^2 \left[ 1 + \log \left( \frac{\mu^2}{m^2} \right) \right] +\mathcal{O}(\varepsilon, m^{-2} )\right\} .  \label{toy EFT loop of fermions} 
\end{align}
%

式\eqref{toy EFT scalar loop diagram}の図は、EFTと完全な理論に同じように寄与するため、両理論の1ループ補正を比較すると消える。したがって、有効作用の1ループ補正は、$\overline{MS}$法での式\eqref{toy EFT loop of fermions}の図と等しい寄与を与えなければならない。
%
\begin{align}
	\parbox[t][8mm][b]{30mm}{\begin{fmfgraph*}(78,50) 
	\fmfleft{i} \fmfright{o} 
	\fmf{dashes}{i,v,o} 
	\fmffreeze
	\fmfv{d.sh=circle,d.f=empty,d.si=.2w,l=$\!\!\!\!\! 1$}{v}
	\end{fmfgraph*}} =  - \frac{i y^2}{4 \pi^2} \left\{\frac{p^2}{2} \log \left( \frac{\mu^2}{m^2} \right) + m^2 \left[ 1 + \log \left( \frac{\mu^2}{m^2} \right) \right] +\mathcal{O}(\varepsilon, m^{-2} )\right\} .
\end{align}
%
この結果から、例えば\ref{sec: matching at loop level}小節で行ったように、$\phi$について2次の$S^{(1)}_{\rm eff}$の形を容易に推定できる。この結果をツリーレベルの結果に加えると、1ループの有効作用の2次部分は、次のようになる。
%
\begin{align}
	S^{(0)}_{\rm eff} + S^{(1)}_{\rm eff} = \int d^4 x \left\{ - \tfrac{1}{2} \mathcal{Z} (\mu) \partial_\mu  \phi \, \partial^\mu \phi - \tfrac{1}{2} \mathcal{M}^2 (\mu) \phi^2 + \dots \right\} \ ,
\end{align}
%
ここで、省略記号は高階微分や$\phi$のより高いべきを含む項を表す。また、Wilson係数を次のように定義した。
%
\begin{subequations}
\begin{align}
	\mathcal{Z} (\mu) &= 1 + \frac{y^2}{8 \pi^2} \log \left( \frac{\mu^2}{m^2} \right)  , \\
	\mathcal{M}^2 (\mu) &= M^2 + \frac{y^2}{4 \pi^2} m^2 \left[ 1 + \log \left( \frac{\mu^2}{m^2} \right) \right] \label{1 loop correction mass scalar} .
\end{align}
\end{subequations}
%

スカラーについて今計算した質量二乗の係数$\mathcal{M}^2 (\mu)$と、スカラーを積分消去してフェルミオンについて得た係数（式\eqref{M 1-loop toy}を参照）の間には、非常に重要な定性的違いがある。フェルミオンの質量パラメータへの1ループ補正は、ツリーレベルのフェルミオン質量$m$に比例していたため、大きな対数がないような$\mu$の値では$\mathcal{M}(\mu) \sim m$となる。言い換えると、フェルミオンのEFTの質量パラメータ$\mathcal{M}(\mu)$の大きさは、UVスケール$M$ではなく、IRスケール$m$によって決まっていた。その結果、結合定数$y$をどの値に固定しても、フェルミオン質量をEFTの強結合スケールより任意に小さく選ぶことができる。

一方、スカラー質量$M$への1ループ補正は、重いフェルミオンの質量$m$に比例する。すなわち、IRスケールではなくUVスケールによって決まる。このことは、$\mathcal{M}^2(\mu)$がその1ループ寄与より小さくはならないと期待すべきことを意味する。すなわち、
%
\begin{align} \label{fine tuning 1}
	\mathcal{M}^2(\mu) \gtrsim \frac{y^2}{4 \pi^2} \, m^2 , 
\end{align}
% 
が$\mu \sim \mathcal{O}(m)$で成り立つと期待される。前の場合と異なり、今度は$y$をどの値に固定しても、$\mathcal{M}^2(m)$に下限がある。ある特定のスケール$\mu$で、ツリーレベルとループの寄与を繊細に\high{微調整}することで、この下限を回避することはできる。しかし、繰り込みスケールを$\mathcal{O}(1)$の相対的な因子だけ変えると、その微調整は保たれないため、これは「不自然な」可能性と見なされる。それに対し「自然な」期待は、スカラー質量が強結合スケール$m$の近傍にあり、そこからの幅の大きさが結合定数$y$で決まるというものである。自然の現実的な理論に適用する際、この期待は、私たちの宇宙の性質はかなり一般的なものであるはずだ、という考えに基づいている。

スカラー質量がUVスケール$m$に敏感であることは、低エネルギーでEFTへ切り替えたことによる見かけの効果ではないと強調しておこう。完全な理論だけを用いても、同様の結果を導ける。$\overline{MS}$繰り込み法では、スカラーの物理的質量は、繰り込まれた質量パラメータ$M$とは一致せず、スカラー伝播関数の極の位置として定義される。すべての1PIループ補正の和を$i \Pi (k^2)$と書くと、物理的質量は$M_{\rm ph}^2 = M^2 - \Pi (-M_{\rm ph}^2)$という関係で決まる。1ループでは、次式を得る。
%
\begin{align}
	M_{\rm ph}^2 = M^2 \! \left\{ 1 - \frac{\lambda}{32 \pi^2} \left[ 1 + \log \left( \frac{\mu^2}{M^2} \right) \right] - \frac{y^2}{8 \pi^2} \log \left( \frac{\mu^2}{M^2} \right) \! \right\} + \frac{y^2}{4 \pi^2} \, m^2 \!\left[ 1 + \log \left( \frac{\mu^2}{m^2} \right) \right] \! , 
\end{align}
%
これから、$M_{\rm ph}^2$を$\frac{y^2}{4 \pi^2} \, m^2 \simeq \frac{y^2}{4 \pi^2} \, m_{\rm ph}^2$よりはるかに小さくするには、微調整が必要だと結論できる。

スカラー質量の\high{UV感度}の問題は、完全な理論のループを鋭いカットオフ$\Lambda$で正則化し、スカラー質量に、カットオフの2乗で発散する補正が生じることに注目して議論されることが多い。一方、フェルミオンの質量は$\Lambda$に対して対数的にしか依存しない。したがってスカラー質量はUV物理により敏感である、というのがその議論である。この主張にはもちろん真実の核がある。しかし学生の頃、私はいつもこれをわかりにくいと感じていた。カットオフ正則化は、ループ計算の中間段階を意味のあるものにするための技巧にすぎないからである。最終的に$\Lambda$の痕跡はすべて、繰り込まれた場と結合定数に吸収されるのだから、発散が2次か対数的かで、なぜ違いが生じるのだろうか。要点は、運動量空間で鋭いカットオフを課すことを、スケール$\Lambda$で働き始めるUV物理の変更と見なせることにある。もっとも、それは非常に非物理的な変更ではある。この見方では、2次発散は再び、スカラー質量の自然な大きさが「UV物理」で決まることの徴候と解釈できる。次元正則化では物理的な質量スケールだけを扱うため、最終的に重要なのはループがどう発散するかではなく、スカラーの物理的質量が物理的なUVスケールにどう依存するかであることが明確になる。
 
説明のつかない微調整がないという要請は、別の仕方でも表せる。明らかに、式\eqref{fine tuning 1}の下限を満たすには、
%
\begin{align} \label{t Hooft naturalness}
	M^2 \gtrsim \frac{y^2}{4 \pi^2} \, m^2 .
\end{align}
%
であれば十分である。式\eqref{1 loop correction mass scalar}によれば、この要請は、UVスケール$\mu = m$において、EFTの質量パラメータの1ループ補正が、そのツリーレベルの値より大きくないことを意味する。この観察から、自然性を少し精密に定義できる。EFTのWilson係数は、その大きさがループ補正によって大きく変わらなければ、技術的な意味で自然、略して技術的に自然であると見なす。このように述べると、\high{技術的自然性}の要請は、Wilson係数が適切に定義された摂動展開をもつことを求めるのに等しい。

式\eqref{full theory}のトイモデルは、$m^2 \ll M^2$なら技術的に自然だが、$M^2 \ll \frac{y^2}{4 \pi^2} \, m^2$ならそうではないことを見た。結局のところ、フェルミオン質量$m$へのループ補正は$m$自身に比例するのに対し、スカラー質量二乗$M^2$へのループ補正は、$M^2$ではなく$m^2$に比例することがすべてである。この違いはどこから来るのだろうか。実は、このトイモデルは$m = 0$のとき、次の離散対称性を追加で獲得する。
%
\begin{align} \label{some sort of chiral symmetry}
	\Psi \to \gamma^5 \Psi \ , \qquad \qquad \qquad \phi \to - \phi \ .
\end{align}
%
量子異常がなければ、ループ補正は理論の対称性を保つはずであり、この模型では量子異常は問題にならない。このため、質量パラメータへの補正がゼロでなくなりうるのは、ツリーレベルですでに$m$がゼロでない場合だけである。言い換えると、$m \to 0$の極限で対称性\eqref{some sort of chiral symmetry}を保つためには、これらの補正もその極限で消えなければならない。より一般に、任意の無次元パラメータ（この場合には$m/M$）をゼロにすると模型の対称性が増すなら、そのパラメータが1よりはるかに小さいことは自然である。この場合、そのパラメータの小ささは、追加の対称性によって「保護されている」という。この形の自然性の基準は、\high{'t Hooft自然性}~\cite{tHooft:1979bh}と呼ばれる。これは、ここで導入した技術的自然性より少し厳しい要請である。何らかの対称性によって保護されていなくても、特定のWilson係数が{\it いかなる}ループ補正も受けない理論が存在するからである。\footnote{'t Hooft自然性と技術的自然性の区別は、文献で一般的ではないものの、有用だと私は考えている。ここで't Hooft自然性と呼んでいるものが、しばしば技術的自然性と呼ばれることに注意してほしい。}とりわけ顕著な例は、Galileon理論~\cite{Goon:2016ihr}である。

この数十年の間に、自然性は、標準模型（SM）を超える物理を探究するうえでの指導原理の一つとなった。この概念が素粒子物理学のコミュニティの研究方針を形作るうえで果たしてきた重要な役割を考えると、そのより広い含意に手短に触れずに済ませるわけにはいかない。ただし、その前に断っておくべきことがある。以下の議論は、この主題についての私個人の見方を反映している。さまざまな資料に当たり、自分自身の見方を育ててほしい。その出発点となる文献を、次の小節にいくつか挙げておいた。

式\eqref{t Hooft naturalness}の自然性の基準は、ここで考えてきた特定のトイモデルに固有のものであり、式\eqref{full theory}の作用のパラメータへの制約である、と強調することが大切だと思う。重いフェルミオンを、別の重い粒子のスペクトルに置き換える、すなわちスカラーのEFTに別の\high{UV完備化}を選べば、この制約は異なる形になる。したがって、スカラー質量が自然かどうかを確かめるには、（{\it i}）スカラーEFTの正しいUV完備化を知り、（{\it ii}）例えば軽いスカラーと重い粒子の両方が関与する散乱過程を調べることで、そのパラメータを実験的に測定する必要がある。このことは、自然性を、低エネルギーEFTがその可能なUV完備化に課す制約と見なすべきことを示唆する。\footnote{これは、UV完備化について何らかの仮定を置くことで低エネルギーEFTに課せる制約とは対照的である。その例として、Lorentz不変なUV完備化をもつ理論におけるスピンと統計の関係~\cite{Weinberg:1995mt}やWilson係数に対する正値性の制約~\cite{Adams:2006sv}、弦理論に由来するEFTの対称性や粒子構成への制約~\cite{Vafa:2005ui}などがある。}

素粒子物理学で最も有名な二つの「自然性の問題」は、一般相対論（GR）に結合したSMの、UVに敏感な二つのパラメータ、すなわちHiggs質量パラメータと宇宙定数に関わる。Higgs質量のUV感度は\high{階層性問題}として知られ、本質的には、スカラーのトイEFTで出会ったものと同じ現象である。階層性問題は、弱い相互作用のスケールを超える物理について、微調整を要しない模型を書くことを、確かに難しくする。それが本当の問題、すなわち自然について私たちが理解していない側面になるかどうかは、SMの正しいUV完備化がどのようなものであるかに依存する。

今のところは、広い視野を保って、弱い相互作用のスケールを超える新しい物理を探すべきである。同時に、素粒子物理学の未解決問題（例えばバリオン数生成、暗黒物質、ニュートリノ質量、強いCP問題など）のいずれも、UVでの解決を必ずしも必要とはしないことを、心に留めておく必要がある。もちろん、そのような解決策の一部は、特に魅力的である。私の考えでは、新しいUV物理を支持する最も強い論拠は、GRがPlanckスケールで有効でなくなる有効理論だという事実である。EFTについての経験から、このスケール以下で新しい自由度が現れ、重力相互作用をより完全に記述すると期待される。また、この新しい物理がHiggs質量パラメータに大きな量子補正を与えると期待するのも妥当である。残念ながら、Planckスケールの重力についての現在の理解では、詳細なマッチング計算を行って、この期待を明示的に確かめることはできない。

\high{宇宙定数問題}は、少なくとも二つの理由から、自然性と素粒子物理学におけるその役割についての理解に、はるかに深刻な難問を突き付ける。第一に、そして何よりも、階層性問題と違って、新しいUV物理の存在に依存しない。質量$m$をもつどんなSM粒子を積分消去しても、宇宙定数には$\Delta \Lambda =\mathcal{O}(1) \frac{m^4}{16 \pi^2 \mpl^2}$という形の1ループ補正が生じる。電子だけでも$\mathcal{O}(10^{-52}) \, \text{GeV}^2$の寄与が生じ、これは観測値$\mathcal{O}(10^{-84}) \, \text{GeV}^2$よりはるかに大きい。第二に、宇宙定数を量子補正から保護できる既知の対称性は、超対称性だけである。残念ながら、これまで超対称性の実験的証拠は得られていない。したがって、超対称性が存在するとしても、弱い相互作用のスケールより高いスケール$M_{\rm susy}$で破れていなければならない。この場合、宇宙定数は一般に$\Delta \Lambda =\mathcal{O}(1) \frac{M_{\rm susy}^4}{16 \pi^2 \mpl^2} \gtrsim \mathcal{O}(10^{-34}) \, \text{GeV}^2$の程度の補正を受け、これも観測値よりはるかに大きい。この状況から、宇宙定数問題の解決は、自然の基本法則における新しい構成原理を示唆するかもしれない、と多くの研究者が考えるようになった。興味深い可能性として、人間原理による考察~\cite{Weinberg:1987dv}や、宇宙論的な緩和機構~\cite{Alberte:2016izw,Graham:2019bfu}がある。
%
\begin{eBox}
%
{\bf 演習2.6：} 規約により、宇宙定数は$S \supset \int \sqrt{-g} \left( - \mpl^2 \Lambda \right)$という形で有効作用に入る。Higgs粒子を積分消去したときに生じる、宇宙定数の1ループのマッチング補正を計算せよ。SMの結合定数の最低次で計算し、結果を自由Higgs場のエネルギー運動量テンソルの真空期待値と比較せよ。
%
\end{eBox}
%



\subsectionlink{さらに学ぶための文献}

EFTの手法については、例えば文献~\cite{Georgi:1994qn,Manohar:1996cq,Pich:1998xt,Rothstein:2003mp,Kaplan:2005es,Burgess:2007pt,Skiba:2010xn,Gripaios:2015qya,Petrov:2016azi,Manohar:2018aog,Cohen:2019wxr}など、優れた総説が数多くある。とりわけ、生成汎関数のレベルで1ループのマッチング計算を実行したわかりやすい例は、文献~\cite{Burgess:2007pt}に見いだせる。マッチングと走りについての非常に教育的な議論は、文献~\cite{Cohen:2019wxr}を参照してほしい。自然性についての興味深い議論は、文献~\cite{Dine:2015xga,Giudice:2017pzm}を参照されたい。一方、階層性問題への挑戦的な見方については、文献~\cite{Manohar:2018aog}の付録を参照してほしい。
宇宙定数問題に関する古典的な文献は、\cite{Weinberg:1988cp}である。


\newpage


